\documentclass[11pt]{article}

\usepackage[margin=1in]{geometry}
\usepackage[T1]{fontenc}
\usepackage{amsmath,amssymb}
\usepackage{booktabs}
\usepackage{xcolor}
\usepackage{hyperref}
\hypersetup{
  colorlinks = true,
  linkcolor  = blue!70!black,
  citecolor  = green!50!black,
  urlcolor   = magenta!70!black
}

\newcommand{\ehtimnote}[1]{%
  \par\smallskip\noindent
  \colorbox{gray!12}{\parbox{\dimexpr\linewidth-2\fboxsep}{\small\textit{In ehtim:} #1}}%
  \par\smallskip}

\title{Interferometric Sampling and Data Corruption\\in the Event Horizon Telescope}
\author{Aaron Barrios}
\date{\today}

\begin{document}
\maketitle

\noindent This note describes how a very-long-baseline interferometer (VLBI)
such as the Event Horizon Telescope (EHT) measures a source, how much of the
image's Fourier information it actually records, and how those measurements
are corrupted. Interstellar scattering and intrinsic source variability are
neglected throughout; the source is treated as a static total-intensity image
$I(x,y)$ on the sky.

Shaded notes point to where each step is implemented in \texttt{ehtim}
(eht-imaging, \url{https://github.com/achael/eht-imaging}) \cite{Chael2018};
names refer to version 1.3.2.

%----------------------------------------------------------------------
\section{What an Interferometer Measures}

\subsection{Visibilities}

Let $(x,y)$ be angular coordinates on the sky (in radians) relative to the
phase center, and let $I(x,y)$ be the sky brightness. Two antennas $i$ and $j$
separated by the baseline vector $\mathbf{b}_{ij}$ record electric fields
$E_i(t)$ and $E_j(t)$. Their time-averaged correlation is, by the
van Cittert--Zernike theorem \cite{TMS2017}, the \emph{complex visibility}
\begin{equation}
V(u,v) = \iint I(x,y)\, e^{-2\pi i\,(ux + vy)}\,dx\,dy,
\label{eq:vis}
\end{equation}
i.e.\ the two-dimensional Fourier transform of the image evaluated at the
spatial frequency $(u,v)$. Here $(u,v)$ is the projection of
$\mathbf{b}_{ij}$ onto the plane perpendicular to the line of sight,
measured in wavelengths $\lambda$. A baseline of projected length $|\mathbf{b}|$
probes angular structure on the scale
\begin{equation}
\theta \sim \frac{\lambda}{|\mathbf{b}|}.
\end{equation}

Because $I$ is real, the visibilities are Hermitian,
\begin{equation}
V(-u,-v) = V^*(u,v),
\end{equation}
so every measurement fixes two conjugate points of the $(u,v)$ plane.

\ehtimnote{the discrete transform is built in
\texttt{observing.obs\_helpers.ftmatrix}. It uses the opposite sign,
$e^{+2\pi i(ux+vy)}$, with $(u,v)$ defined from site~1 minus site~2, and its
$x$ axis increases toward the east (pixel columns run in decreasing $x$). It
also multiplies each visibility by the Fourier transform of a pixel ``pulse''
(triangle by default). A custom pipeline must match these conventions to agree
with \texttt{ehtim}.}

\subsection{Baseline coordinates and Earth-rotation synthesis}

Write the baseline in an equatorial frame with components
$(L_X, L_Y, L_Z)$, where $Z$ points to the celestial pole, $X$ toward hour
angle $H = 0$ on the celestial equator, and $Y$ toward $H = -6^{\rm h}$. For a
source at declination $\delta$ observed at hour angle $H$ \cite{TMS2017},
\begin{align}
u &= \frac{1}{\lambda}\left(L_X\sin H + L_Y\cos H\right), \\
v &= \frac{1}{\lambda}\left(-L_X\sin\delta\cos H + L_Y\sin\delta\sin H + L_Z\cos\delta\right).
\end{align}
As the Earth rotates, $H$ advances and each baseline traces an arc of the
ellipse
\begin{equation}
u^2 + \left(\frac{v - v_0}{\sin\delta}\right)^2 = \frac{L_X^2 + L_Y^2}{\lambda^2},
\qquad v_0 = \frac{L_Z\cos\delta}{\lambda},
\end{equation}
together with its Hermitian reflection. An array of $N$ stations has
\begin{equation}
N_{\rm bl} = \frac{N(N-1)}{2}
\end{equation}
baselines, each contributing one such track while the source is above the
horizon at both stations.

\ehtimnote{station positions and SEFDs are read from an array file with
\texttt{eh.array.load\_txt} (example files such as
\texttt{arrays/EHT2017\_m87.txt} are in the GitHub repository).
\texttt{Image.observe(array, tint, tadv, tstart, tstop, bw, ...)} or
\texttt{Array.obsdata(...)} generates the tracks, with elevation limits
\texttt{elevmin}~$=10^\circ$ and \texttt{elevmax}~$=85^\circ$ by default.
\texttt{Obsdata.plotall('u','v')} plots the coverage.}

\subsection{Resolution and scale}

The EHT observes at $\nu \approx 230$~GHz ($\lambda \approx 1.3$~mm) with
baselines up to $|\mathbf{b}|_{\max} \approx 10{,}700$~km, so
\begin{equation}
|\mathbf{u}|_{\max} = \frac{|\mathbf{b}|_{\max}}{\lambda} \approx 8\times10^{9}\,\lambda
\quad\Longrightarrow\quad
\theta_{\min} \approx \frac{\lambda}{|\mathbf{b}|_{\max}} \approx 25~\mu{\rm as}.
\end{equation}
For comparison, the angular gravitational radius
$\theta_g = GM/(c^2 D)$ is $\approx 3.8~\mu$as for M87$^*$ \cite{EHT2019I} and
$\approx 4.8~\mu$as for Sgr~A$^*$ \cite{EHT2022I}, so the nominal resolution
is $\sim 5$--$7\,\theta_g$. The $n \ge 1$ photon rings, whose widths are a
small fraction of $\theta_g$, are therefore unresolved.

%----------------------------------------------------------------------
\section{Incomplete Sampling of Fourier Space}

\subsection{The sampling function}

An observation measures the visibility only at a finite set of points
$\{(u_k, v_k)\}_{k=1}^{K}$, one per baseline per integration time. Define the
sampling function
\begin{equation}
S(u,v) = \sum_{k=1}^{K}\left[\delta(u - u_k)\,\delta(v - v_k) + \delta(u + u_k)\,\delta(v + v_k)\right].
\end{equation}
The measured data are equivalent to $S(u,v)\,V(u,v)$.

\subsection{Dirty image and dirty beam}

Inverting the sampled visibilities directly gives the \emph{dirty image},
\begin{equation}
I_D(x,y) = \iint S(u,v)\,W(u,v)\,V(u,v)\,e^{2\pi i\,(ux + vy)}\,du\,dv
= \left(I * B_D\right)(x,y),
\label{eq:dirty}
\end{equation}
where $W(u,v)$ is a chosen weighting (natural weighting $W = 1/\sigma^2$,
uniform weighting $W \propto 1/\text{local sample density}$), $*$ denotes
convolution, and the \emph{dirty beam} is
\begin{equation}
B_D(x,y) = \iint S(u,v)\,W(u,v)\,e^{2\pi i\,(ux + vy)}\,du\,dv.
\end{equation}
The dirty beam is the interferometer's point-spread function. Unlike a
telescope PSF it has strong sidelobes and negative regions, determined
entirely by the geometry of the $(u,v)$ coverage.

\ehtimnote{\texttt{Obsdata.dirtyimage(npix, fov, weighting=...)} and
\texttt{Obsdata.dirtybeam(npix, fov, weighting=...)}, with
\texttt{weighting} set to \texttt{'uniform'} (default) or \texttt{'natural'}.}

\subsection{What is lost}

Equation~\eqref{eq:dirty} shows the essential difference from ordinary
blurring. A Gaussian PSF multiplies the Fourier transform by a smooth,
everywhere-positive function: high spatial frequencies are attenuated but
still present. The sampling function $S$ is \emph{zero} almost everywhere.
Any image component $\Delta I$ whose Fourier transform vanishes on all sampled
points,
\begin{equation}
\widetilde{\Delta I}(u_k, v_k) = 0 \quad \forall k,
\end{equation}
produces identical data and is invisible to the instrument (the null space of
the measurement). Reconstructing such components requires prior assumptions
about the image, not data.

Coverage is incomplete for several reasons:
\begin{itemize}
  \item \textbf{Few stations.} The 2017 campaign used eight stations at six
        geographic sites \cite{EHT2019II}, giving at most a few tens of
        distinct tracks.
  \item \textbf{No baselines beyond an Earth diameter.} Spatial frequencies
        above $|\mathbf{u}|_{\max}$ are never measured, which sets the
        resolution limit above.
  \item \textbf{Few short baselines.} Only intra-site pairs (e.g.\ ALMA--APEX,
        SMA--JCMT) probe large angular scales; structure between those scales
        and the shortest inter-site baseline is poorly constrained.
  \item \textbf{Limited mutual visibility.} A baseline exists only while the
        source is above the elevation limit at both stations, so tracks are
        short arcs rather than full ellipses.
  \item \textbf{Anisotropic coverage.} The distribution of stations makes
        the coverage, and hence the effective beam, elongated; resolution
        differs between directions on the sky.
  \item \textbf{Lost data.} Weather, technical failures, and scans dropped
        during calibration remove additional points.
\end{itemize}

\ehtimnote{coverage can be degraded by removing stations
(\texttt{Array.make\_subarray}, \texttt{Obsdata.flag\_sites}) or time ranges
(\texttt{Obsdata.flag\_UT\_range}), or by shortening \texttt{tstart}--\texttt{tstop}
in \texttt{observe}.}

\subsection{Restoring beam}

For display, reconstructed images are often convolved with a
\emph{restoring} (``CLEAN'') beam: an elliptical Gaussian fitted to the main
lobe of $B_D$ \cite{Hogbom1974}. For the EHT this has a full width at half
maximum of $\sim 20~\mu$as \cite{EHT2019IV}. The restoring beam is a
presentation choice that suppresses unconstrained high-frequency structure;
it is not the physical measurement process.

\ehtimnote{\texttt{Obsdata.fit\_beam()} returns the fitted
$(\text{FWHM}_{\rm maj}, \text{FWHM}_{\rm min}, \theta)$, and
\texttt{Image.blur\_gauss(beamparams)} convolves an image with it.}

%----------------------------------------------------------------------
\section{Corruption of the Measured Visibilities}

\subsection{Measurement equation}

Neglecting polarization, the visibility recorded on baseline $ij$ at time $t$
is
\begin{equation}
V_{ij}^{\rm obs}(t) = g_i(t)\,g_j^*(t)\,V_{ij}(t) + \epsilon_{ij}(t),
\label{eq:rime}
\end{equation}
where $V_{ij}(t) = V(u_{ij}(t), v_{ij}(t))$ is the true visibility,
$g_i = |g_i|\,e^{i\phi_i}$ is a complex \emph{station gain}, and
$\epsilon_{ij}$ is thermal noise.

\subsection{Thermal noise}

The noise $\epsilon_{ij}$ is complex Gaussian with zero mean and standard
deviation (per real and imaginary component)
\begin{equation}
\sigma_{ij} = \frac{1}{\eta_Q}\sqrt{\frac{{\rm SEFD}_i\,{\rm SEFD}_j}{2\,\Delta\nu\,\tau}},
\qquad
{\rm SEFD} = \frac{2k_B\,T_{\rm sys}^*}{\eta_A\,A_{\rm geom}},
\label{eq:thermal}
\end{equation}
where $\Delta\nu$ is the bandwidth, $\tau$ the integration time,
$\eta_Q \approx 0.88$ the quantization efficiency for two-bit sampling,
$T_{\rm sys}^*$ the system temperature corrected for atmospheric opacity,
$\eta_A$ the aperture efficiency, and $A_{\rm geom}$ the collecting area
\cite{TMS2017}. The system equivalent flux density (SEFD) ranges from
$\sim 10^2$~Jy for phased ALMA to $\sim 10^4$~Jy for the smaller single
dishes, so baselines to ALMA are by far the most sensitive. Noise in the
visibility domain maps to spatially correlated noise in the image, shaped by
the dirty beam; it is not independent from pixel to pixel.

\ehtimnote{Eq.~\eqref{eq:thermal} is \texttt{observing.obs\_helpers.blnoise},
including the $1/0.88$ factor; it is applied when
\texttt{observe(..., add\_th\_noise=True)} (the default), using the SEFDs in
the array file.}

\subsection{Station phase errors}

At millimeter wavelengths, turbulent water vapor in the troposphere changes
the signal path length above each station on timescales of seconds. The
resulting phases $\phi_i(t)$ are effectively random, which limits coherent
integration to $\tau \sim 10$~s and makes the absolute phase of
$V_{ij}^{\rm obs}$ unusable without calibration.

\subsection{Station amplitude errors}

Uncertainties in $T_{\rm sys}^*$, aperture efficiency, pointing, and focus
give gain amplitudes $|g_i|$ that typically deviate from unity by
$\sim 5$--$20\%$ and can vary from scan to scan.

\ehtimnote{in \texttt{observe}, \texttt{phasecal=False} adds random station
phases and \texttt{ampcal=False} adds random gain amplitudes, with fractional
scatter \texttt{gainp} and systematic offset \texttt{gain\_offset} (both
$0.1$ by default). \texttt{stabilize\_scan\_phase} and
\texttt{stabilize\_scan\_amp} hold the errors constant within a scan.}

\subsection{Non-closing systematics}

Errors that do not factor into station-based terms (e.g.\ baseline-dependent
effects, residual polarization leakage) are commonly represented by an
additional fractional noise floor added in quadrature,
\begin{equation}
\sigma_{ij}^2 \;\to\; \sigma_{ij}^2 + \left(f_{\rm sys}\,|V_{ij}|\right)^2,
\end{equation}
with $f_{\rm sys}$ of order a percent \cite{EHT2019IV}.

\ehtimnote{\texttt{Obsdata.add\_fractional\_noise(f\_sys)}.}

%----------------------------------------------------------------------
\section{Closure Quantities}

Because the gains in Eq.~\eqref{eq:rime} are attached to stations, certain
combinations of visibilities are independent of them.

\paragraph{Closure phase.} For a triangle of stations $(i,j,k)$,
\begin{equation}
\psi_{ijk} = \arg V_{ij}^{\rm obs} + \arg V_{jk}^{\rm obs} + \arg V_{ki}^{\rm obs},
\end{equation}
in which every station phase $\phi_i$ cancels.

\paragraph{Closure amplitude.} For a quadrangle $(i,j,k,l)$,
\begin{equation}
A_{ijkl} = \frac{|V_{ij}^{\rm obs}|\,|V_{kl}^{\rm obs}|}{|V_{ik}^{\rm obs}|\,|V_{jl}^{\rm obs}|},
\end{equation}
in which every gain amplitude $|g_i|$ cancels.

\paragraph{Information content.} With $N$ stations observing simultaneously,
the numbers of independent closure phases and closure amplitudes are
\begin{equation}
N_\psi = \frac{(N-1)(N-2)}{2}, \qquad N_A = \frac{N(N-3)}{2},
\end{equation}
compared with $N(N-1)/2$ visibility phases and amplitudes. Closure phases
discard the absolute position of the image; closure amplitudes discard the
total flux. At high signal-to-noise, the noise in each closure quantity is
approximately
\begin{equation}
\sigma_\psi \approx \sqrt{\sum_{\rm baselines} \left(\frac{\sigma}{|V|}\right)^2},
\qquad
\sigma_{\ln A} \approx \sqrt{\sum_{\rm baselines} \left(\frac{\sigma}{|V|}\right)^2},
\end{equation}
summed over the three or four baselines involved \cite{Blackburn2020}.

\ehtimnote{\texttt{Obsdata.c\_phases(count='min')} and
\texttt{Obsdata.c\_amplitudes(count='min', ctype='logcamp')} return
independent sets; \texttt{count='max'} returns all (redundant) combinations.}

%----------------------------------------------------------------------
\section{Image Reconstruction}

Recovering $I(x,y)$ from the corrupted, incomplete data is an ill-posed
inverse problem: infinitely many images fit the data equally well. Two
families of methods are used by the EHT \cite{EHT2019IV}.

\paragraph{CLEAN.} The image is modeled as a sum of point sources found
iteratively by subtracting scaled, shifted copies of $B_D$ from $I_D$, with
self-calibration of the gains \cite{Hogbom1974}.

\paragraph{Regularized maximum likelihood (RML).} The image is found by
minimizing
\begin{equation}
J(I) = \sum_{\rm data} \alpha_d\,\chi^2_d(I) + \sum_R \beta_R\,\mathcal{R}(I),
\end{equation}
where the $\chi^2_d$ terms compare the model with visibility amplitudes,
closure phases, or closure amplitudes, and the regularizers $\mathcal{R}$
(maximum entropy, total variation, sparsity, etc.)\ encode prior assumptions
about the image \cite{Chael2016,Chael2018}. A learned reconstruction network
plays the role of the regularizer: it supplies the prior that fills the null
space of the measurement.

\ehtimnote{RML imaging is \texttt{Imager(obs, init\_im, data\_term=\{...\},
reg\_term=\{...\})} followed by \texttt{make\_image\_I()}. Data terms include
\texttt{'vis'}, \texttt{'amp'}, \texttt{'cphase'} and \texttt{'logcamp'};
regularizers include \texttt{'simple'} (maximum entropy), \texttt{'tv'},
\texttt{'tv2'}, \texttt{'l1w'} and \texttt{'gs'}. Self-calibration of the
gains is \texttt{calibrating.self\_cal.self\_cal}; CLEAN-style routines are in
\texttt{imaging/clean.py}, though dedicated packages (e.g.\ DIFMAP) are
standard for CLEAN.}

%----------------------------------------------------------------------
\section{Connecting Simulated Images to Visibilities}

A ray-traced image computed on a grid in units of $GM/c^2$ must be converted
to angular units before Eq.~\eqref{eq:vis} can be applied:
\begin{equation}
x\,[{\rm rad}] = \theta_g\,\alpha, \qquad
\theta_g = \frac{GM}{c^2 D}.
\end{equation}
For a pixel size $\Delta x$ and field of view $L$, the discrete Fourier
transform samples spatial frequencies up to $1/(2\Delta x)$ with spacing
$1/L$. Avoiding aliasing of the longest EHT baselines requires
\begin{equation}
\Delta x < \frac{1}{2|\mathbf{u}|_{\max}} \approx 12~\mu{\rm as},
\end{equation}
which is easily satisfied by ray-traced images ($\Delta x \sim 1~\mu$as), and
the field of view must be large enough that the image goes to zero at its
edges. Because the sampled $(u_k, v_k)$ do not lie on the FFT grid,
visibilities are evaluated by direct or non-uniform Fourier transforms
\cite{Chael2018}.

\ehtimnote{a ray-traced array becomes
\texttt{eh.image.Image(image, psize, ra, dec, rf=...)} with \texttt{psize}
in radians and pixel values in Jy. The transform in \texttt{observe} is
selected by \texttt{ttype}: \texttt{'dtft'} (direct, exact but slow),
\texttt{'nfft'} (default; needs the \texttt{pynfft} package) or
\texttt{'fast'} (padded FFT plus interpolation, padding set by
\texttt{fft\_pad\_factor}).}

%----------------------------------------------------------------------
\begin{thebibliography}{9}

\bibitem{TMS2017}
A.~R. Thompson, J.~M. Moran, and G.~W. Swenson,
\emph{Interferometry and Synthesis in Radio Astronomy}, 3rd ed.\
(Springer, Cham, 2017), doi:10.1007/978-3-319-44431-4 (open access).

\bibitem{EHT2019I}
Event Horizon Telescope Collaboration,
``First M87 Event Horizon Telescope results. I. The shadow of the
supermassive black hole,''
\emph{Astrophys.\ J.\ Lett.} \textbf{875}, L1 (2019), arXiv:1906.11238.

\bibitem{EHT2019II}
Event Horizon Telescope Collaboration,
``First M87 Event Horizon Telescope results. II. Array and instrumentation,''
\emph{Astrophys.\ J.\ Lett.} \textbf{875}, L2 (2019), arXiv:1906.11239.

\bibitem{EHT2019IV}
Event Horizon Telescope Collaboration,
``First M87 Event Horizon Telescope results. IV. Imaging the central
supermassive black hole,''
\emph{Astrophys.\ J.\ Lett.} \textbf{875}, L4 (2019), arXiv:1906.11241.

\bibitem{EHT2022I}
Event Horizon Telescope Collaboration,
``First Sagittarius A$^*$ Event Horizon Telescope results. I. The shadow of
the supermassive black hole in the center of the Milky Way,''
\emph{Astrophys.\ J.\ Lett.} \textbf{930}, L12 (2022),
doi:10.3847/2041-8213/ac6674.

\bibitem{Hogbom1974}
J.~A. H\"ogbom,
``Aperture synthesis with a non-regular distribution of interferometer
baselines,''
\emph{Astron.\ Astrophys.\ Suppl.\ Ser.} \textbf{15}, 417 (1974).

\bibitem{Chael2016}
A.~A. Chael \emph{et al.},
``High-resolution linear polarimetric imaging for the Event Horizon
Telescope,''
\emph{Astrophys.\ J.} \textbf{829}, 11 (2016), arXiv:1605.06156.

\bibitem{Chael2018}
A.~A. Chael \emph{et al.},
``Interferometric imaging directly with closure phases and closure
amplitudes,''
\emph{Astrophys.\ J.} \textbf{857}, 23 (2018), arXiv:1803.07088.

\bibitem{Blackburn2020}
L.~Blackburn \emph{et al.},
``Closure statistics in interferometric data,''
\emph{Astrophys.\ J.} \textbf{894}, 31 (2020), arXiv:1910.02062.

\end{thebibliography}

\end{document}
